\documentclass[12pt,a4paper]{article}
\usepackage[utf8]{inputenc}
\usepackage[brazil]{babel}
\usepackage{graphicx}
\usepackage{amssymb, amsfonts, amsmath}
\usepackage{float}
\usepackage{enumerate}
\usepackage[top=2.5cm, bottom=2.5cm, left=1.25cm, right=1.25cm]{geometry}

\begin{document}
\pagestyle{empty}

\begin{center}
  \begin{tabular}{ccc}
    \begin{tabular}{c}
      \includegraphics[scale=0.25]{../../biblioteca/imagem/brasao-de-armas-brasil} \\
    \end{tabular} & 
    \begin{tabular}{c}
      Ministério da Educação \\
      Universidade Federal dos Vales do Jequitinhonha e Mucuri \\
      Faculdade de Ciências Sociais, Aplicadas e Exatas - FACSAE \\
      Departamento de Ciências Exatas - DCEX \\
      Disciplina: Introdução à Lógica \quad Semestre: 2022/1\\
      Prof. Dr. Luiz C. M. de Aquino\\
      Aluno(a):\rule{6cm}{0.1mm} \quad Data: \rule{0.5cm}{0.1mm}/\rule{0.5cm}{0.1mm}/\rule{1cm}{0.1mm}\\
    \end{tabular} &
    \begin{tabular}{c}
      \includegraphics[scale=0.25]{../../biblioteca/imagem/logo-ufvjm} \\
    \end{tabular}
  \end{tabular}
\end{center}

\begin{center}
 \textbf{Avaliação I}
\end{center}

\textbf{Instruções}
\begin{itemize}
 \item Todas as justificativas necessárias na solução de cada questão devem 
 estar presentes nesta avaliação;
 \item As respostas finais de cada questão devem estar escritas de caneta;
 \item Esta avaliação tem um total de 30,0 pontos.
\end{itemize}

\begin{enumerate}
  \item \textbf{[6,0 pontos]} Complete a tabela verdade abaixo.
  
  \begin{center}
    \begin{tabular}{|c|c|c|c|c|c|c|c|c|c|}
      \hline
      $p$ & $q$ & $r$ & $\lnot p$ & $\lnot q$ & 
      $\lnot q$ & $(\lnot r) \rightarrow (\lnot q)$ & 
      $(\lnot p)\leftrightarrow q$ & 
      $((\lnot r) \rightarrow (\lnot q))\vee ((\lnot p)\leftrightarrow q)$\\ \hline
       & & & & & & & & \\ \hline
       & & & & & & & &  \\ \hline
       & & & & & & & &  \\ \hline
       & & & & & & & &  \\ \hline
       & & & & & & & &  \\ \hline
       & & & & & & & &  \\ \hline
       & & & & & & & &  \\ \hline
       & & & & & & & &  \\ \hline
    \end{tabular}
  \end{center}

  \item \textbf{[6,0 pontos]} Complete a tabela verdade abaixo.
  
  \begin{center}
    \begin{tabular}{|c|c|c|c|c|c|c|}
      \hline
      $p$ & $q$ & $r$ &
      $p\wedge q \wedge r$ &
      $p\vee q \vee r$ &
      $p\to (q\to r)$ &
      $p\leftrightarrow (q\leftrightarrow r)$\\ \hline
       & & & & & & \\ \hline
       & & & & & & \\ \hline
       & & & & & & \\ \hline
       & & & & & & \\ \hline
       & & & & & & \\ \hline
       & & & & & & \\ \hline
       & & & & & & \\ \hline
       & & & & & & \\ \hline
    \end{tabular}
  \end{center}
    
  \item \textbf{[6,0 pontos]} Considerando a proposição condicional ``se eu 
  gosto de maça, então eu não gosto de banana'', escreva as suas respectivas 
  proposições:

    \begin{enumerate}
      \item recíproca;

      \rule{\linewidth}{0.1mm}
      
      \rule{\linewidth}{0.1mm}
            
      \item contrária;

      \rule{\linewidth}{0.1mm}
      
      \rule{\linewidth}{0.1mm}

      \item contrapositiva;

      \rule{\linewidth}{0.1mm}
      
      \rule{\linewidth}{0.1mm}
      
      \item negação.

      \rule{\linewidth}{0.1mm}
      
      \rule{\linewidth}{0.1mm}
    \end{enumerate}

  \item \textbf{[6,0 pontos]} Escreva a negação de cada proposição composta abaixo:

    \begin{enumerate}
      \item Eu não gosto de sorvete de baunilha e eu gosto de doce de batata;

      \rule{\linewidth}{0.1mm}
      
      \rule{\linewidth}{0.1mm}
            
      \item Eu tenho 20 anos ou eu não estou solteiro;

      \rule{\linewidth}{0.1mm}
      
      \rule{\linewidth}{0.1mm}

      \item Eu faço pão se, e somente se, eu faço manteiga;

      \rule{\linewidth}{0.1mm}
      
      \rule{\linewidth}{0.1mm}
      
      \item Se eu não vou cozinhar feijão, então eu vou cozinhar beterraba.

      \rule{\linewidth}{0.1mm}
      
      \rule{\linewidth}{0.1mm}
    \end{enumerate}

    \item \textbf{[6,0 pontos]} (Vunesp/TJ-SP - Adaptada) Em um edifício com apartamentos 
    somente nos andares de 1º ao 4º, moram 4 meninas, 
    em andares distintos: Ana, Fran, Alice e Maria, não necessariamente nessa ordem. 
    Cada uma delas tem um animal de estimação diferente: gato, cachorro, passarinho e hamster, 
    não necessariamente nessa ordem. Fran vive reclamando do barulho feito pelo passarinho, 
    no andar imediatamente acima do seu. Maria, que não mora no 4º andar, mora um andar acima do de Ana, 
    que tem o gato e não mora no 2º andar. Quem mora no 3º andar tem um hamster. 
    Com base nessas informações, complete a tabela abaixo com Verdadeiro ou Falso para 
    determinar em qual andar mora cada menina e qual é seu respectivo animal de estimação.
    
  
  \begin{center}
    \begin{tabular}{|c|c|c|c|c|c|c|c|c|}
      \hline
       Nome & 1º & 2º & 3º & 4º & Cachorro & Gato & Passarinho & Hamster \\ \hline
       Ana & & & & & & & & \\ \hline
       Fran & & & & & & & & \\ \hline
       Alice & & & & & & & & \\ \hline
       Maria & & & & & & & & \\ \hline
    \end{tabular}
  \end{center}

\end{enumerate}

\end{document}